\documentclass[12pt,a4paper]{article}
\usepackage{graphicx}
\usepackage{amsmath}
\usepackage{hyperref}
\usepackage{xcolor}
\usepackage{geometry}
\geometry{margin=1in}

\title{\textbf{Smart Calorie \& Nutrition Analyzer}}
\author{Usman Khan \\ IMSC Data Science \\ ISA Project 2025}
\date{}

\begin{document}
\maketitle

\section*{1. Problem Statement}
In today's world, maintaining a healthy lifestyle is a growing challenge due to busy schedules, poor dietary habits, and limited understanding of nutritional requirements. 
Most people consume food without knowing its calorie composition or daily energy expenditure, leading to weight gain or nutritional deficiencies.  
To address this issue, the project \textbf{Smart Calorie \& Nutrition Analyzer} aims to build an intelligent tool that helps users:
\begin{itemize}
    \item Calculate personal daily calorie requirements based on individual attributes.
    \item Look up the nutritional values of common foods.
    \item Predict caloric values using machine learning from macronutrient data.
\end{itemize}
This application enables data-driven dietary decisions and promotes healthier eating habits.

\section*{2. Solution}
The Smart Calorie \& Nutrition Analyzer is implemented using Python and the Streamlit framework for interactive web applications. 
It integrates user input, nutritional datasets, and a predictive model to offer three main functionalities:

\begin{enumerate}
    \item \textbf{Personal Calorie Calculator:} Computes Basal Metabolic Rate (BMR) and Total Daily Energy Expenditure (TDEE) using user-provided data such as weight, height, age, and activity level.
    \item \textbf{Food Nutrition Lookup:} Allows users to search for food items from a dataset (\texttt{food\_data.csv}) and view their macronutrient and calorie content.
    \item \textbf{Calorie Prediction Model:} Employs Linear Regression to estimate calorie content based on protein, carbohydrate, and fat values.
\end{enumerate}

The application also offers user interface customization with Dark and Light modes, improving usability and accessibility.

\section*{3. Visualization}
To provide meaningful insights, the system uses several visualization features that assist users in understanding their dietary composition:
\begin{itemize}
    \item \textbf{Pie Chart:} Displays the percentage distribution of protein, carbohydrates, and fat in a given input.
    \item \textbf{Bar Chart:} Compares macronutrient grams for visual inspection of diet balance.
    \item \textbf{Interactive Table:} Displays food search results dynamically using \texttt{st.dataframe()}, allowing users to explore food data effectively.
\end{itemize}

All visualizations are done using \texttt{Matplotlib} and \texttt{Seaborn} libraries.

\section*{4. Formulae}

\subsection*{4.1 Basal Metabolic Rate (BMR)}
For calculating the basal metabolic rate, the Mifflin-St Jeor equation is used:
\[
\text{BMR}_{male} = 10W + 6.25H - 5A + 5
\]
\[
\text{BMR}_{female} = 10W + 6.25H - 5A - 161
\]
where:
\[
W = \text{weight (kg)}, \quad H = \text{height (cm)}, \quad A = \text{age (years)}
\]

\subsection*{4.2 Total Daily Energy Expenditure (TDEE)}
The TDEE is computed using:
\[
\text{TDEE} = \text{BMR} \times \text{Activity Factor}
\]

Activity factors used in this project:
\begin{align*}
\text{Sedentary:} & \quad 1.2 \\
\text{Lightly Active:} & \quad 1.375 \\
\text{Moderately Active:} & \quad 1.55 \\
\text{Active:} & \quad 1.725 \\
\text{Very Active:} & \quad 1.9
\end{align*}

\subsection*{4.3 Calorie Prediction Model}
The calorie prediction model is based on the Linear Regression equation:
\[
\hat{y} = b_0 + b_1x_1 + b_2x_2 + b_3x_3
\]
where:
\[
\hat{y} = \text{predicted calories}, \quad x_1 = \text{protein (g)}, \quad x_2 = \text{carbs (g)}, \quad x_3 = \text{fat (g)}
\]

\section*{5. Conclusion}
The Smart Calorie \& Nutrition Analyzer serves as an intelligent, interactive, and data-driven diet management tool. 
By combining user health metrics, machine learning, and real-time visualization, the project offers an accessible way to understand and manage personal nutrition effectively.

\end{document}
